\subsubsection{Classical Renormalization of Codimension-two Brane Couplings --- arXiv:0710.4598}

\textbf{Authors:} Claudia de Rham

\paragraph{主な主張}
codimension-two 以上のbraneでは局在matter sourceに伴う特異性が現れるが、brane couplingを古典的にくりこむことで、低エネルギー理論をUV regularizationの取り方に依存しない形で定義できることを示した\cite{DeRham:2007ClassicalBrane}。

\paragraph{新規性}
brane上で評価したbulk propagatorから生じるUV感度を量子loopではなくclassical RG flowとして整理し、codimension-two brane EFTにおける局在結合の系統的なくりこみを提示した。

\paragraph{位置付け}
codimension-two brane相互作用の一般的なEFT基盤を与える論文であり、2次元bulk-neutrino模型で現れる対数的KK cutoff依存をbrane couplingのrunningとして解釈する際の基礎となる。
